% from ega4-18-fr.tex
\oldpage[IV-4]{160}

\begin{remark}[18.10.3.1]
\label{IV.18.10.3.1}
Si $f:X\to Y$ est localement de présentation finie, dominant et quasi-compact, la fibre $f^{-1}(\eta)$ est non vide en vertu de (I, 1.1.5). Par contre, si l'on ne suppose pas $f$ quasi-compact, $f$ peut être non ramifié et dominant sans que $f$ soit étale (en supposant toujours vérifiées les hypothèses de (18.10.3) pour $X$ et $Y$). C'est ce que montre l'exemple suivant: on prend pour $Y$ le plan affine $\Spec(\mathbf C[T,U])$; on considère d'autre part une famille $(X_j)_{j\in\mathbf Z}$ de préschémas isomorphes à la droite affine $\Spec(\mathbf C[T])$, et on forme le préschéma obtenu en « recollant » $X_{2j}$ et $X_{2j+1}$ au point $-1$ et $X_{2j+1}$ et $X_{2j+2}$ au point $+1$. Définissons d'autre part $f:X\to Y$ comme étant égal sur $X_{2j}$ à l'immersion fermée dont l'image est la droite d'équation $y=2j$, transformant le point $-1$ en $(2j-1,2j)$ et le point $+1$ en $(2j+1,2j)$; sur $X_{2j+1}$, $f_j$ est l'immersion fermée dont l'image est la droite d'équation $x=2j+1$, transformant le point $-1$ en $(2j+1,2j)$ et le point $-1$ en $(2j+1,2j+2)$. Il est clair que $f$ est une immersion locale (donc est non ramifié) et est dominant mais la fibre du point générique de $Y$ est vide et $f$ n'est pas étale.
\end{remark}

\begin{corollary}[18.10.4]
\label{IV.18.10.4}
Soient $g:Y\to S$, $h:X\to S$ deux morphismes localement de présentation finie, $f:X\to Y$ un $S$-morphisme, $x$ un point de $X$, $y=f(x)$, $s=g(y)=h(x)$. On suppose que $h$ est plat au point $x$ et $g^{-1}(s)$ normal au point $y$. Alors, pour que $f$ soit étale au point $x$, il faut et il suffit que $f$ soit non ramifié au point $x$ et que $\dim_xh^{-1}(s)=\dim_yg^{-1}(s)$.
\end{corollary}

Les conditions sont nécessaires en vertu du fait que si $f$ est étale au point $x$, il en est de même de $f_s:h^{-1}(s)\to g^{-1}(s)$, et la conclusion résulte alors de ce que $f_s$ est plat et quasi-fini au point $x$ (6.1.2). Inversement, si les conditions de l'énoncé sont vérifiées, pour voir que $f$ est étale au point $x$, il suffit, en vertu de (17.8.3) de voir que $f_s$ l'est. On peut donc se borner au cas où $S$ est le spectre d'un corps. Il résulte alors de l'hypothèse $\dim_xX=\dim_yY$ et du fait que $f$ est non ramifié au point $x$ que l'on a (par (5.2.3) et (17.4.1)) $\dim\sh{O}_{X,x}=\dim\sh{O}_{Y,y}$; puisque $f$ est quasi-fini au point $x$ et $\sh{O}_{Y,y}$ intègre, on déduit de ($0_{\mathrm I}$, 7.4.4) et ($0$, 16.3.10) que l'homomorphisme $\sh{O}_{Y,y}\to\sh{O}_{X,x}$ est injectif; il suffit alors pour conclure d'appliquer (18.10.1).

\begin{corollary}[18.10.5]
\label{IV.18.10.5}
Avec les notations de (18.10.4), on suppose que $h$ est plat et que $g^{-1}(s)$ est normal pour tout $s\in S$ (ce qui sera par exemple le cas si $g$ est lisse). Pour que $f$ soit une immersion ouverte, il faut et il suffit que $f$ soit un monomorphisme et que, pour tout $x\in X$, on ait $\dim_xh^{-1}(s)=\dim_yg^{-1}(s)$, où $y=f(x)$ et $s=g(y)=h(x)$.
\end{corollary}

On n'a à prouver que la suffisance des conditions; or, il résulte de (17.1.3) qu'un monomorphisme localement de présentation finie est non ramifié, donc les hypothèses entraînent que $f$ est étale (18.10.4); on conclut alors par (17.9.1).

\begin{lemma}[18.10.6]
\label{IV.18.10.6}
Soit $f:X\to Y$ un morphisme plat et localement quasi-fini ($\mathrm{Err}_{\mathrm{III}}$, 20).
\begin{enumerate}
  \item[(i)] L'ensemble $\operatorname{Max}(X)$ des points maximaux de $X$ est en correspondance biunivoque canonique avec l'ensemble $\coprod_{y\in\operatorname{Max}(Y)}f^{-1}(y)$.
  \item[(ii)] Si $Y$ est irréductible et de point générique $\eta$, l'ensemble des composantes irréductibles de $X$ est en correspondance biunivoque canonique avec $f^{-1}(\eta)$, et en particulier est fini si et seulement si $f^{-1}(\eta)$ est fini.

